\documentclass[
    12pt,
    a4paper,
    oneside,
    brazil,
    sumario=tradicional
]{abntex2}

\usepackage[utf8]{inputenc}
\usepackage[T1]{fontenc}
\usepackage{times}
\usepackage{microtype}
\usepackage{indentfirst}
\usepackage{xcolor}
\usepackage{amsmath,amssymb,amsfonts}
\usepackage{listings}
\usepackage{booktabs}
\usepackage{tabularx}
\usepackage{hyperref}

\usepackage[
    top=3cm, bottom=2cm,
    left=3cm, right=2cm
]{geometry}

\hypersetup{
    colorlinks=true,
    linkcolor=blue,
    urlcolor=blue,
    citecolor=blue,
    pdftitle={RepairDAO},
    pdfauthor={Arthur Henrique de Souza Paro}
}

\lstset{
    basicstyle=\ttfamily\small,
    breaklines=true,
    frame=single,
    backgroundcolor=\color{gray!10},
    keywordstyle=\color{blue}\bfseries,
    commentstyle=\color{gray},
    stringstyle=\color{orange},
    numbers=left,
    numberstyle=\tiny\color{gray},
    numbersep=5pt,
    tabsize=2,
    showstringspaces=false
}

% ---------------------------------------------------------------
% Informações do trabalho
% ---------------------------------------------------------------
\titulo{RepairDAO: Um Protocolo Web3 Descentralizado para Marketplace de\\
Serviços Técnicos com Reputação Portável via NFT}
\autor{Arthur Henrique de Souza Paro}
\local{Brasil}
\data{Abril de 2026}
\instituicao{Residência em TIC 29 --- Unidade 1, Capítulo 5 \par
             Curso de Desenvolvimento Web 3.0}
\tipotrabalho{Relatório Técnico}
\preambulo{Relatório técnico desenvolvido para o Curso de Desenvolvimento
Web 3.0 da Residência em TIC 29, Unidade 1, Capítulo 5.}

\begin{document}

\frenchspacing

% ---------------------------------------------------------------
% Elementos pré-textuais
% ---------------------------------------------------------------
\pretextual

\imprimircapa

\imprimirfolhaderosto

% Resumo
\begin{resumo}
Este trabalho apresenta o \textbf{RepairDAO}, um protocolo descentralizado
desenvolvido em Solidity para a rede Ethereum que implementa um marketplace
de serviços técnicos (manutenção e reparo de eletrônicos) com reputação
portável via NFT, pagamento automatizado em escrow, staking com recompensas
dinâmicas ajustadas por oráculo de preço Chainlink e governança comunitária
via DAO. O protocolo é composto por quatro contratos inteligentes:
\texttt{RepairToken} (ERC-20), \texttt{RepairNFT} (ERC-721 soulbound),
\texttt{RepairEscrow} e \texttt{RepairDAO}. Todos foram deployados e testados
na testnet Sepolia, com auditoria de segurança realizada via Slither 0.11.5.

\vspace{\onelineskip}
\noindent\textbf{Palavras-chave}: blockchain. Solidity. ERC-20. ERC-721. NFT.
escrow. DAO. Chainlink. staking. Web3.
\end{resumo}

\tableofcontents*

\cleardoublepage

% ---------------------------------------------------------------
% Elementos textuais
% ---------------------------------------------------------------
\textual

% -------------------------------------------------------------------
\section{Introdução}

Plataformas centralizadas de serviços técnicos, como GetNinjas,
apresentam três problemas estruturais: (1)~o técnico paga para receber o
contato do cliente sem garantia de pagamento; (2)~a reputação acumulada
fica presa na plataforma e é perdida em caso de bloqueio de conta;
(3)~taxas e regras são definidas unilateralmente pelo operador.

O \textbf{RepairDAO} resolve esses problemas através de contratos
inteligentes na blockchain Ethereum:

\begin{itemize}
    \item \textbf{Reputação como NFT portável}: o histórico de
    serviços e avaliações do técnico é armazenado em um NFT
    ERC-721 intransferível (\textit{soulbound}), pertencente ao técnico,
    não à plataforma.

    \item \textbf{Pagamento em escrow automático}: os tokens ficam
    bloqueados no contrato até a confirmação do cliente, eliminando o
    risco de inadimplência.

    \item \textbf{Taxa de protocolo justa}: 2\% por transação vão ao
    tesouro da DAO, gerida pelos próprios participantes.

    \item \textbf{Recompensas para stakers}: quem deposita tokens em staking
    recebe recompensas dinâmicas ajustadas pelo preço do ETH via
    oráculo Chainlink.

    \item \textbf{Governança descentralizada}: alterações de
    parâmetros e resolução de disputas são decididas por
    votação ponderada pelo saldo em staking.
\end{itemize}

O projeto foi desenvolvido como MVP prático para o Curso de Web~3.0
da Residência em TIC~29 e deployado na testnet Sepolia para validação
funcional completa.

% -------------------------------------------------------------------
\section{Arquitetura do Protocolo}

O protocolo é composto por quatro contratos que se comunicam via interfaces
Solidity. As dependências externas são OpenZeppelin~v5 e
\texttt{@chainlink/contracts}.

\subsection{RepairToken (ERC-20)}

O \texttt{RepairToken} é o token de utilidade do protocolo, símbolo
\texttt{REPAIR}, supply máximo de 10~milhões de unidades. Herda de
\texttt{ERC20} e \texttt{Ownable} da OpenZeppelin.

Características principais:
\begin{itemize}
    \item Supply máximo: $10^7 \times 10^{18}$~wei, verificado no
    constructor e na função \texttt{mintReward}.
    \item Mint restrito: apenas o \texttt{RepairEscrow} pode mintar tokens
    adicionais (recompensas de staking), controlado pela variável
    \texttt{escrowContract} configurada via \texttt{setEscrowContract}, que
    só pode ser chamada uma vez pelo owner.
    \item Supply inicial mintado para o owner no deploy, parametrizado via
    constructor.
\end{itemize}

\begin{lstlisting}[language=Solidity, caption={Mint restrito ao Escrow}]
function mintReward(address to, uint256 amount)
    external
{
    require(msg.sender == escrowContract,
        "apenas o escrow pode mintar");
    require(totalSupply() + amount <= MAX_SUPPLY,
        "acima do supply maximo");
    _mint(to, amount);
}
\end{lstlisting}

\subsection{RepairNFT (ERC-721 Soulbound)}

NFT de reputação intransferível --- um por técnico, emitido
automaticamente no primeiro serviço concluído. A
intransferibilidade é garantida sobrescrevendo \texttt{transferFrom}
e \texttt{safeTransferFrom} para sempre reverter:

\begin{lstlisting}[language=Solidity, caption={Bloqueio de transferência}]
function transferFrom(address, address, uint256)
    public pure override
{
    revert("NFT de reputacao e intransferivel");
}
\end{lstlisting}

A estrutura \texttt{Reputacao} armazenada em cada token contém:
\begin{itemize}
    \item \texttt{totalServicos}: número acumulado de serviços.
    \item \texttt{somaAvaliacoes}: soma das notas (1--5).
    \item \texttt{dataPrimeiroServico}: timestamp on-chain.
\end{itemize}

A média de avaliações é calculada como:
$$\bar{A} = \frac{\displaystyle\sum_{i=1}^{n} a_i}{n}, \quad a_i \in \{1,2,3,4,5\}$$

\subsection{RepairEscrow}

Contrato central do protocolo. Gerencia o fluxo de serviço, staking e
a integração com o oráculo.

\subsubsection{Fluxo de Serviço}

\[
\texttt{Aberto}
  \xrightarrow{\texttt{confirmarServico}}  \texttt{Concluido}
\]
\[
\texttt{Aberto}
  \xrightarrow{\texttt{abrirDisputa}}      \texttt{Disputado}
  \xrightarrow{\texttt{resolverDisputa}}   \texttt{Concluido}
\]

\begin{enumerate}
    \item \textbf{criarServico}: cliente deposita tokens REPAIR no contrato
    (escrow). O valor equivalente em USD é calculado via oráculo no momento
    da criação.

    \item \textbf{confirmarServico}: cliente confirma a entrega e avalia o
    técnico (1--5). O contrato distribui: 98\% ao técnico, 2\% ao tesouro
    da DAO. Em seguida chama \texttt{registrarServico} no NFT.

    \item \textbf{abrirDisputa}: coloca o serviço no estado
    \texttt{Disputado} para votação na DAO.

    \item \textbf{resolverDisputa}: owner (DAO) direciona os fundos ao
    vencedor após deliberação.
\end{enumerate}

\subsubsection{Staking e Recompensas}

A recompensa acumulada por um staker é calculada pela fórmula:

\begin{equation}
R = \frac{V \cdot \tau \cdot t \cdot M}{D \cdot T_a \cdot 100}
\end{equation}

onde:
\begin{itemize}
    \item $V$ = valor em staking (wei)
    \item $\tau$ = taxa base anual em BPS (1\,000 = 10\%)
    \item $t$ = tempo em staking (segundos)
    \item $M$ = multiplicador: 150 se ETH $>$ \$2\,000, senão 100
    \item $D$ = denominador BPS (10\,000)
    \item $T_a$ = 1 ano em segundos ($365 \times 86\,400$)
\end{itemize}

\subsection{RepairDAO}

Governança comunitária com dois tipos de proposta:

\begin{itemize}
    \item \textbf{AlterarParametro}: proposta de mudança de parâmetro
    do protocolo.
    \item \textbf{ResolverDisputa}: proposta de resolução de disputa
    de um serviço específico.
\end{itemize}

Parâmetros de governança:
\begin{itemize}
    \item Duração de votação: 3 dias.
    \item Mínimo para propor: 1\,000 REPAIR em staking.
    \item Peso do voto: proporcional ao saldo em staking.
\end{itemize}

Ciclo completo de uma proposta:
\[
\texttt{Ativa}
\xrightarrow{\texttt{finalizarProposta}}
\texttt{Aprovada} \mid \texttt{Rejeitada}
\xrightarrow{\texttt{executarProposta}}
\texttt{Executada}
\]

% -------------------------------------------------------------------
\section{Segurança}

\subsection{Medidas Implementadas}

\begin{table}[h]
\caption{Mecanismos de segurança por contrato}
\label{tab:seguranca}
\begin{tabularx}{\textwidth}{@{}lX@{}}
\toprule
\textbf{Mecanismo} & \textbf{Aplicação} \\
\midrule
\texttt{ReentrancyGuard}      & RepairEscrow, RepairDAO \\
\texttt{Ownable}              & Todos os contratos \\
Modifier \texttt{apenasEscrow}& RepairToken, RepairNFT \\
\texttt{setEscrowContract} único & RepairToken, RepairNFT \\
NFT intransferível            & RepairNFT \\
\texttt{pragma \^{}0.8.20}   & Overflow/underflow nativo \\
\texttt{require} em entradas  & Todas as funções públicas \\
\bottomrule
\end{tabularx}
\fonte{Elaborado pelo autor (2026).}
\end{table}

\subsection{Proteção contra Reentrada}

As funções críticas (\texttt{confirmarServico}, \texttt{retirar},
\texttt{resolverDisputa}) são protegidas pelo modifier
\texttt{nonReentrant} da OpenZeppelin, seguindo o padrão
\textit{checks-effects-interactions}: o estado é alterado antes das
transferências externas.

\subsection{Controle de Acesso}

O protocolo utiliza controle de acesso em dois níveis:
\begin{enumerate}
    \item \textbf{Owner}: funções administrativas e configuração
    inicial.
    \item \textbf{escrowContract}: mint de tokens e atualização do
    NFT --- configurado uma única vez, imutável após a
    configuração.
\end{enumerate}

\subsection{Auditoria com Slither}

Auditoria est\'atica realizada com \textbf{Slither~0.11.5} e compilador
\texttt{solc~0.8.28}:

\begin{lstlisting}[language=bash, caption={Comando de auditoria}]
py -3.10 -m slither contracts/ \
  --solc-remaps \
  "@openzeppelin=node_modules/@openzeppelin \
   @chainlink=node_modules/@chainlink"
\end{lstlisting}

\textbf{Resultado}: \textbf{zero vulnerabilidades críticas} nos contratos
do protocolo. Os alertas gerados referem-se exclusivamente a contratos da
OpenZeppelin e Chainlink (dependências externas já auditadas).

% -------------------------------------------------------------------
\section{Integração com Oráculo Chainlink}

\subsection{Feed ETH/USD na Sepolia}

O \texttt{RepairEscrow} integra o feed de preço \texttt{ETH/USD} da
Chainlink na testnet Sepolia:

\begin{center}
\small\texttt{0x694AA1769357215DE4FAC081bf1f309aDC325306}
\end{center}

A leitura \'e realizada via \texttt{AggregatorV3Interface}:

\begin{lstlisting}[language=Solidity, caption={Leitura do or\'aculo}]
function getPrecoETH() public view returns (uint256) {
    (, int256 price,,,) =
        priceFeed.latestRoundData();
    require(price > 0, "preco invalido");
    return uint256(price);
}
\end{lstlisting}

O valor retornado possui 8 casas decimais (padrão Chainlink):
$$P_{\$} = \frac{\texttt{getPrecoETH()}}{10^8}$$

\subsection{Multiplicador Dinâmico de Recompensa}

Se $P_{\$} > 2\,000$, o multiplicador é $1{,}5\times$; caso contrário,
$1{,}0\times$:

\begin{lstlisting}[language=Solidity, caption={L\'ogica do multiplicador}]
// ETH > $2000: precoETH > 200_000_000_000
uint256 multiplicador =
    precoETH > 200_000_000_000 ? 150 : 100;
\end{lstlisting}

\subsection{Resultado Obtido na Sepolia}

Teste realizado em 27 de abril de 2026:
\begin{itemize}
    \item \texttt{getPrecoETH()} retornou: \texttt{228\,999\,000\,000}
    \item Preço ETH em USD: \textbf{\$2\,289,99}
    \item Multiplicador ativo: \textbf{1,5$\times$}
\end{itemize}

% -------------------------------------------------------------------
\section{Scripts de Interação (Etapa~5)}

Três scripts Node.js foram desenvolvidos utilizando \textbf{ethers.js~v6}
e \textbf{Hardhat~2.22}, executados via
\texttt{npx hardhat run --network sepolia}.

\subsection{scripts/stake.js}

Fluxo: \texttt{approve} $\to$ \texttt{depositar} $\to$ consulta
recompensa pendente + preço ETH via oráculo.

\textbf{Evidências} (TXs confirmadas na Sepolia):
\begin{itemize}
    \item Approve: \texttt{0x...cc07c2}
    \item Depositar: \texttt{0x...9fbce}
\end{itemize}
\textbf{Resultado}: 1\,000~REPAIR em staking, pool total = 1\,000~REPAIR.

\subsection{scripts/mintNFT.js}

Fluxo: \texttt{approve} $\to$ \texttt{criarServico} $\to$
\texttt{confirmarServico(1, 5)} $\to$ verificação do NFT.

\textbf{Evidências} (TXs confirmadas na Sepolia):
\begin{itemize}
    \item criarServico: \texttt{0x9590fef0...}
    \item confirmarServico: \texttt{0x22da641d...}
\end{itemize}
\textbf{Resultado}: NFT Token~\#1 mintado para o técnico,
média~=~\textbf{5/5}.

\subsection{scripts/votar.js}

Fluxo: verifica poder de voto $\to$
\texttt{proporAlteracaoParametro} $\to$ \texttt{votar(1, true)}.

\textbf{Evidências} (TXs confirmadas na Sepolia):
\begin{itemize}
    \item proporAlteracaoParametro: \texttt{0x93579937...}
    \item votar: \texttt{0x57233da2...}
\end{itemize}
\textbf{Resultado}: Proposta~\#1 criada, \textbf{1\,000~REPAIR} a favor,
expiração em 30/04/2026.

% -------------------------------------------------------------------
\section{Deploy na Testnet Sepolia}

\subsection{Stack Tecnológica}

\begin{table}[h]
\caption{Ambiente de desenvolvimento}
\label{tab:stack}
\begin{tabularx}{\textwidth}{@{}lX@{}}
\toprule
\textbf{Tecnologia} & \textbf{Versão / Detalhe} \\
\midrule
Solidity       & \texttt{\^{}0.8.20}, evmVersion cancun \\
OpenZeppelin   & v5 (ERC20, ERC721, Ownable, ReentrancyGuard) \\
Chainlink      & \texttt{@chainlink/contracts} \\
Hardhat        & 2.22 (CommonJS) \\
ethers.js      & v6 \\
Node.js        & v25.8.0 \\
Slither        & 0.11.5 \\
IDE de Deploy  & Remix IDE \\
Rede           & Sepolia Testnet (Infura) \\
\bottomrule
\end{tabularx}
\fonte{Elaborado pelo autor (2026).}
\end{table}

\subsection{Endereços dos Contratos Deployados}

\begin{table}[h]
\caption{Contratos deployados na Sepolia}
\label{tab:deploy}
\begin{tabularx}{\textwidth}{@{}lX@{}}
\toprule
\textbf{Contrato} & \textbf{Endereço (Sepolia)} \\
\midrule
RepairToken  & \texttt{0xc57AFe0D023c016a00fAAaC1F4E406...ff7} \\
RepairNFT    & \texttt{0xF40F843aabB16f7781be1090594A15...bc7} \\
RepairEscrow & \texttt{0xD0c9f72a98d28eF6259541e4cDb8cc...b55} \\
RepairDAO    & \texttt{0x44ad016590bf6873E2cf89eA98bd14...eC}  \\
\bottomrule
\end{tabularx}
\fonte{Elaborado pelo autor (2026).}
\end{table}

Links completos no Etherscan Sepolia:
\begin{itemize}
    \item \url{https://sepolia.etherscan.io/address/0xc57AFe0D023c016a00fAAaC1F4E4068b39891ff7}
    \item \url{https://sepolia.etherscan.io/address/0xF40F843aabB16f7781be1090594A157C32A10bc7}
    \item \url{https://sepolia.etherscan.io/address/0xD0c9f72a98d28eF6259541e4cDb8ccc2Fb29ab55}
    \item \url{https://sepolia.etherscan.io/address/0x44ad016590bf6873E2cf89eA98bd1424cF7854eC}
\end{itemize}

\subsection{Ordem de Deploy e Configurações Pós-Deploy}

\begin{enumerate}
    \item \texttt{RepairToken(owner, 1\,000\,000)}
    \item \texttt{RepairNFT(owner)}
    \item \texttt{RepairEscrow(owner, token, nft, chainlinkFeed, tesouro)}
    \item \texttt{RepairDAO(owner, token, escrow)}
    \item \texttt{RepairToken.setEscrowContract(escrowAddress)}
    \item \texttt{RepairNFT.setEscrowContract(escrowAddress)}
\end{enumerate}

% -------------------------------------------------------------------
\section{Resultados e Discussão}

\subsection{Funcionalidades Validadas}

Todas as funcionalidades do MVP foram validadas na Sepolia:

\begin{itemize}
    \item \textbf{Oráculo}: \texttt{getPrecoETH()} retornou \$2\,289,99
    com multiplicador 1,5$\times$ ativo.
    \item \textbf{Staking}: depósito de 1\,000~REPAIR confirmado on-chain.
    \item \textbf{Escrow}: criação e confirmação de serviço
    com mint automático do NFT de reputação.
    \item \textbf{NFT}: Token~\#1 mintado com média 5/5 armazenada on-chain.
    \item \textbf{Governança}: proposta criada e voto registrado com
    1\,000~REPAIR de peso.
\end{itemize}

\subsection{Limitações do MVP}

\begin{itemize}
    \item O preço de referência para REPAIR é aproximado via feed
    ETH/USD --- em produção exigiria feed ou AMM dedicados.
    \item \texttt{executarProposta} é restrita ao owner, centralizando a
    execução; em produção seria substituída por execução
    automática pós-aprovação.
    \item Não há mecanismo de \textit{timelock} entre aprovação e
    execução de propostas.
\end{itemize}

% -------------------------------------------------------------------
\section{Conclusão}

O \textbf{RepairDAO} demonstrou a viabilidade de um marketplace
descentralizado de serviços técnicos, implementando em quatro
contratos inteligentes os mecanismos que eliminam os pontos centrais de
falha das plataformas convencionais: reputação portável via NFT
soulbound, pagamento sem intermediários via escrow automático,
recompensas dinâmicas ajustadas por oráculo de mercado e governança
comunitária.

O protocolo foi validado end-to-end na testnet Sepolia, com todas as
transações auditáveis publicamente no Etherscan. A auditoria estática
com Slither não identificou vulnerabilidades críticas. Os scripts
ethers.js demonstraram a interação programática cobrindo os fluxos
de staking, serviço com mint de NFT e votação na DAO.

Como evoluções possíveis destacam-se: feed REPAIR/USD dedicado,
timelock na execução de propostas, sistema de penalização para
avaliações fraudulentas e migração para mainnet com auditoria formal.

% -------------------------------------------------------------------
\postextual

\begin{thebibliography}{00}

\bibitem{erc20}
VOGELSTELLER, F.; BUTERIN, V\@.
\textbf{EIP-20: Token Standard}.
Ethereum Improvement Proposals, 2015.
Disponível em: \url{https://eips.ethereum.org/EIPS/eip-20}.
Acesso em: 27 abr. 2026.

\bibitem{erc721}
ENTRIKEN, W\@. et al\@.
\textbf{EIP-721: Non-Fungible Token Standard}.
Ethereum Improvement Proposals, 2018.
Disponível em: \url{https://eips.ethereum.org/EIPS/eip-721}.
Acesso em: 27 abr. 2026.

\bibitem{chainlink}
CHAINLINK LABS.
\textbf{Chainlink Data Feeds Documentation}.
2024.
Disponível em: \url{https://docs.chain.link/data-feeds}.
Acesso em: 27 abr. 2026.

\bibitem{openzeppelin}
OPENZEPPELIN.
\textbf{OpenZeppelin Contracts v5.0}.
2024.
Disponível em: \url{https://docs.openzeppelin.com/contracts/5.x/}.
Acesso em: 27 abr. 2026.

\bibitem{slither}
TRAIL OF BITS.
\textbf{Slither: Static Analyzer for Solidity}.
2023.
Disponível em: \url{https://github.com/crytic/slither}.
Acesso em: 27 abr. 2026.

\bibitem{hardhat}
NOMIC FOUNDATION.
\textbf{Hardhat: Ethereum Development Environment}.
2024.
Disponível em: \url{https://hardhat.org}.
Acesso em: 27 abr. 2026.

\bibitem{soulbound}
BUTERIN, V\@.; WEYL, E\@. G\@.; OHLHAVER, P\@.
\textbf{Decentralized Society: Finding Web3's Soul}.
2022.
Disponível em: \url{https://papers.ssrn.com/sol3/papers.cfm?abstract_id=4105763}.
Acesso em: 27 abr. 2026.

\end{thebibliography}

\end{document}
